\section*{Introduction}

A graph $G$ is \emph{$H$-free} if no induced subgraph of $G$ is isomorphic to $H$. Write $\hom(G)$ for
the largest size of a clique or a stable set of $G$. Every graph on $n$ vertices has $\hom(G)\ge
\frac12\log_2 n$, and random graphs show that this is best possible up to the constant. In 1977 Erd\H{o}s
and Hajnal~\cite{EH77,EH} conjectured that excluding a single induced subgraph changes the picture completely.

\begin{conjecture}
For every graph $H$ there exists $c>0$ such that every $H$-free graph $G$ satisfies $\hom(G)\ge|G|^{c}$.
\end{conjecture}

A graph $H$ for which this holds is said to have the \emph{Erd\H{o}s--Hajnal property}. The property is
known for every graph with at most four vertices and is preserved by substitution~\cite{APS}. The bull
was settled by Chudnovsky and Safra~\cite{CS}, and the five-cycle by Chudnovsky, Scott, Seymour and
Spirkl~\cite{C5}. Nguyen, Scott and Seymour proved it for the five-vertex path~\cite{NSS7}, which settled
every graph on five vertices. Since then it has been proved for a few six-vertex graphs: the E-graph and
the bird~\cite{HJZ}, and the graph with edges $ab,bc,cd,de,af,bf,df$~\cite{TN}. For paths in general,
Nguyen, Scott and Seymour~\cite{NSS5} showed that $P_k$-free graphs satisfy $\hom(G)\ge
2^{(\log|G|)^{1-o(1)}}$, and Bousquet, Lagoutte and Thomass\'e~\cite{BLT} proved a polynomial bound when
both $P_k$ and its complement are excluded. After $P_5$, the six-vertex path $P_6$ is the smallest path
for which the conjecture was open.

\begin{theoremL}\label{thm:main}
There exists $\tau>0$ such that every $P_6$-free graph $G$ satisfies $\hom(G)\ge|G|^{\tau}$.
\end{theoremL}

A graph is $P_6$-free if and only if its complement is $\cP$-free, and $\hom(G)=\hom(\Gbar)$, so
Theorem~\ref{thm:main} is equivalent to the Erd\H{o}s--Hajnal property of $\cP$. We work with
$\cP$-free graphs throughout. We in fact prove the stronger \emph{polynomial R\"odl property}
(Theorem~\ref{thm:polyrodl}): there is $A$ such that for every $\eps\in(0,\frac12)$, every $\cP$-free
graph $G$ has an induced subgraph on at least $\eps^{3A}|G|$ vertices whose maximum degree, or whose
complement's maximum degree, is at most $\eps$ times its number of vertices.

The proof determines $\tau$ by a formula (Table~\ref{tab:constants}) in the constant $\delta$ of R\"odl's
theorem and the exponent of the Erd\H{o}s--Hajnal theorem for $P_5$. No numerical values of these are
available, so we cannot give $\tau$ numerically; the formula only shows that $\tau<10^{-15}$.

\subsection*{What is new and what is inherited}
The proof uses four results from the literature as black boxes: R\"odl's theorem, a blockade theorem of Nguyen,
Scott and Seymour for sparse graphs without a long induced path, the Erd\H{o}s--Hajnal theorem for
$P_5$, and the comb lemma of Chudnovsky, Scott, Seymour and Spirkl. The architecture of the proof
(iterative sparsification organised around combs, niceness and a layout theorem) is that of the proof for
$P_5$ by Nguyen, Scott and Seymour~\cite{NSS7}. That proof uses the absence of an induced \emph{house},
the complement of $P_5$, in exactly two places, and both fail for $\cP$-free graphs. The following parts
are new:
\begin{itemize}[leftmargin=1.5em]
\item the three six-vertex certificates: the module law (Lemma~\ref{lem:module}) and the diagonal and
  house lemmas (Lemmas~\ref{lem:diagonal} and~\ref{lem:house});
\item the Tooth Lemma (Lemma~\ref{lem:tooth}), which combines the module law with the modular
  decomposition of a comb's teeth;
\item the resulting comb lemma for $\cP$ (Lemma~\ref{lem:comb}), whose output blockades are only
  \emph{semisparse}, and the observation that the layout theorem accepts semisparse input
  (Theorem~\ref{thm:layout});
\item the house-free-pattern argument that brings the Erd\H{o}s--Hajnal theorem for $P_5$ into the second
  sparsification round (Corollary~\ref{cor:housefree} and Lemma~\ref{lem:round2step}).
\end{itemize}
The remaining steps, namely the two sparsification rounds and the reduction to the polynomial R\"odl
property, follow~\cite{NSS7} with the constants recomputed. We prove them in full as well, in the form in
which they were formalized.

\subsection*{The imported theorems}
The proof uses the following four theorems and nothing else. Section~\ref{sec:imported} quotes each
from its source, next to the Lean proposition that renders it. Theorems I, III and IV come from refereed
papers. Theorem II comes from~\cite{NSS5}, which has not yet appeared in a journal; its proof takes half a
page, and the formalization proves it, so the formal result assumes only Theorems I, III and IV.

\medskip
\noindent\begin{tabularx}{\textwidth}{@{}lXl@{}}
\toprule
& imported theorem & Lean proposition \\
\midrule
I & R\"odl's theorem~\cite{Rodl}, as stated in~\cite[Theorem 1.3]{NSS7}, for $H=\cP$ & \lean{RodlCoP6} \\
II & sparse graphs without an induced path $P$ contain long anticomplete blockades~\cite[statement 3.1]{NSS5}, for $P=P_6$ & \lean{NssPath6} \\
III & the Erd\H{o}s--Hajnal property of $P_5$~\cite[Theorem 1.2]{NSS7} & \lean{EhP5} \\
IV & the comb lemma~\cite[Lemma 4.3]{NSS7}, from~\cite{C5} & \lean{NssComb} \\
\bottomrule
\end{tabularx}

\subsection*{Formal verification}
The whole argument is formalized in Lean~4~\cite{lean4} with Mathlib~\cite{mathlib}. The formal
theorem \lean{EHP6.erdos\_hajnal\_P6\_of\_imported} states that the four propositions above imply
Theorem~\ref{thm:main}. Theorem II is proved in the formalization, so Theorems I, III and IV alone imply
Theorem~\ref{thm:main}; this is the formal theorem \lean{erdos\_hajnal\_P6\_of\_cited}. The proofs of both use no axioms beyond Lean's three standard ones. It was checked with
Lean~4.34.1, replayed from an empty environment with \lean{leanchecker}, and checked by the comparator
tool with two kernels, Lean's own and nanoda, an independent type checker written in Rust. The source and
build instructions are in the repository \repo. Section~\ref{sec:formal} describes the definitions, the
correspondence between the paper and the Lean files, and how to reproduce the checks. The paper can be
read without any knowledge of Lean. Independently of this work, Bonnet has formalized in Lean the
Erd\H{o}s--Hajnal theorems for $C_5$ and $P_5$, including R\"odl's theorem and a version of the comb
lemma~\cite{Lax54,Lax57}. Connecting these to our formalization would turn Theorems I and III, and,
after a change of constant, IV into proved statements; we have not done this.

\subsection*{Priority}
To the best of our knowledge, Theorem~\ref{thm:main} is new. On 27 September 2026 we checked every paper
indexed by Semantic Scholar that cites~\cite{NSS5}, \cite{NSS7}, \cite{C5} or~\cite{HJZ}, searched arXiv
and the web, and read the page and discussion of the conjecture on the Erd\H{o}s problems website (Problem
61), Paul Seymour's list of papers, and the lists of Erd\H{o}s problems resolved with the help of AI. We
found no proof of the Erd\H{o}s--Hajnal property for $P_6$. The closest recent results are~\cite{HJZ}
and~\cite{TN}. By the count of Tran and Nguyen~\cite{TN}, eight complementary pairs of prime six-vertex
graphs were open as of 26 August 2026, and $\{P_6,\cP\}$ is one of them. A public repository by Jared
Wilder~\cite{Wilder} develops local structure for $P_6$-free graphs towards the same goal and states that
it does not prove the full result. We did not have access to MathSciNet, zbMATH or Google Scholar, and we
cannot rule out unpublished work.

\subsection*{Organization}
The next section explains the idea of the proof and where $\cP$-freeness is used.
Section~\ref{sec:prelim} fixes notation, states the four imported theorems, and proves the elementary
tools. Sections~\ref{sec:local}--\ref{sec:comb} contain the new structural part: the local lemmas, the
Tooth Lemma and the comb lemma. Sections~\ref{sec:round1}--\ref{sec:round2} run the two sparsification
rounds. Sections~\ref{sec:conclusion} and~\ref{sec:c01} derive Theorem~\ref{thm:main}, and
Section~\ref{sec:formal} documents the formalization.
